\subtopichead{প্রসঙ্গ, টোকেন ও উত্তর}
একজন শিক্ষার্থী ছবি না-দেখানো HTML কোডের কারণ ও সংশোধন জানতে চাইল। প্রশ্নের সঙ্গে কোড, স্ক্রিনশট, আগের আলোচনা ও উত্তরের নির্দেশও প্রসঙ্গের অংশ হতে পারে। উত্তর হতে পারে সম্ভাব্য কারণ, সংশোধিত কোড ও পরীক্ষা করার ধাপ। শুধু কোড পড়ে মডেল ছবির ফাইল সত্যিই আছে কি না নিশ্চিত করতে পারে না; ফাইল দেখার অনুমোদিত টুল বা পরীক্ষার তথ্য দরকার।
\chOnePair{\chOneNative{LLMDetail}}{প্রম্পট ও প্রসঙ্গ থেকে টোকেন এবং তার সংখ্যারূপ তৈরি হয়। Transformer টোকেনগুলোর সম্পর্ক প্রক্রিয়া করে। পরবর্তী টোকেনের স্কোর থেকে Decoding টোকেন বেছে নেয়। নতুন টোকেনসহ প্রসঙ্গ নিয়ে পরবর্তী ধাপ চলে; থামার শর্ত পূরণ হলে উত্তর সম্পূর্ণ হয়। এটি প্রচলিত Autoregressive Text Generation-এর সরল রূপ।}
টোকেন পূর্ণ শব্দ, শব্দের অংশ, যতিচিহ্ন বা অন্য ছোট একক হতে পারে। বাংলা ও ইংরেজিতে টোকেনের ভাগ একই হতে হবে না। ছবি ও শব্দের ইনপুটেও উপযুক্ত সংখ্যারূপ লাগে। Attention প্রাসঙ্গিক অংশগুলোর সম্পর্কের গুরুত্ব নির্ণয়ে সাহায্য করে। যেমন, ``রাফি সাকিবকে বই দিল; সে বইটি পড়ল'' বাক্যে সর্বনামের প্রসঙ্গ বুঝতে হয়; অস্পষ্টতায় মডেল ভুলও করতে পারে।
\term{Context Window} হলো একবার প্রক্রিয়াকরণে মডেলের তথ্য নেওয়ার সীমা। প্রশ্ন, নথি, কথোপকথন ও তৈরি টোকেন এই সীমার অংশ হতে পারে। দীর্ঘ ফাইল দেওয়া মানেই প্রতিটি তথ্য সমানভাবে ব্যবহৃত হবে না। প্রসঙ্গ, মডেল, টুলের ফল ও টোকেন নির্বাচনের পদ্ধতি বদলালে একই প্রশ্নের উত্তর বদলাতে পারে। কোনো পদ্ধতিতে সম্ভাবনা অনুযায়ী Sampling হয়; সর্বোচ্চ স্কোরের টোকেনই প্রতি ধাপে নেওয়া বাধ্যতামূলক নয়।
\subtopichead{শেখানো, নতুন তথ্য ও মূল্যায়ন}
\begin{compactpoints}
\item \term{Pre-training}: বড় ডেটা থেকে সম্পর্ক শেখানো এবং ভুলের মাপকাঠি অনুযায়ী ওজন সমন্বয় করা হয়। ভাষা মডেলে পরবর্তী টোকেন অনুমান পরিচিত প্রশিক্ষণকাজ।
\item \term{Post-training}: নির্বাচিত উদাহরণ, মানুষের মূল্যায়ন বা অন্য প্রতিক্রিয়া দিয়ে নির্দেশ অনুসরণ ও উত্তরের আচরণ উন্নত করা যায়। পদ্ধতি মডেলভেদে বদলায়।
\item ব্যবহারকালে দেওয়া তথ্য বর্তমান প্রসঙ্গে কাজে লাগতে পারে; তা স্থায়ী প্রশিক্ষণ হওয়ার প্রমাণ নয়। ডেটা রাখা বা ভবিষ্যৎ প্রশিক্ষণে ব্যবহার করা সেবার পৃথক নীতির বিষয়।
\item মূল্যায়নে নতুন উদাহরণ দিয়ে ভাষাগত মান, নির্ভুলতা, ভুল ও কাজের উপযোগিতা পরীক্ষা করতে হয়। শেখানোর উদাহরণ মুখস্থ থাকা যথেষ্ট নয়।
\end{compactpoints}
\term{Retrieval-Augmented Generation বা RAG}-এ প্রশ্ন অনুযায়ী বই বা অনুমোদিত তথ্যভান্ডারের প্রাসঙ্গিক অংশ উদ্ধার করে মডেলের প্রসঙ্গে দেওয়া হয়। মডেল সেই তথ্য ব্যবহার করে উত্তর তৈরি করে। উদ্ধার ভুল হলে উত্তরও ভুল হতে পারে; উত্তর ও উদ্ধৃতি মূল অংশের সঙ্গে মিলিয়ে দেখা দরকার।
Search ও RAG ব্যবহারের সময় বাইরের তথ্য উদ্ধার করে। \term{Fine-tuning}-এ নির্বাচিত প্রশিক্ষণ-ডেটা দিয়ে শেখানো ওজন সমন্বয় করা হয়। নতুন ফাইল যোগ করলেই Fine-tuning হয় না।
\subtopichead{ইনপুট, ফল ও যাচাই}
\begin{compacttable}
\begin{longtable}{|p{\dimexpr.32\linewidth-7pt\relax}|p{\dimexpr.68\linewidth-7pt\relax}|}
\hline\textbf{ইনপুট} & \textbf{ফল ও যাচাই}\\\hline\endfirsthead
\hline\textbf{ইনপুট} & \textbf{ফল ও যাচাই}\\\hline\endhead
প্রশ্ন ও পাঠ্যাংশ & ব্যাখ্যা বা সারসংক্ষেপের সঙ্গে মূল অর্থ মেলাতে হয়।\\\hline
পাতার ছবি ও প্রশ্ন & সম্ভাব্য রোগের লেবেল ছবি ও বিশেষজ্ঞের তথ্য দিয়ে যাচাই করতে হয়।\\\hline
কথার অডিও & লেখা তৈরি হলে নাম, সংখ্যা ও উচ্চারণের ভুল পরীক্ষা করতে হয়।\\\hline
লেখা ও কণ্ঠের নির্দেশ & তৈরি অডিওতে উচ্চারণ, অর্থ ও কণ্ঠ ব্যবহারের অনুমতি দেখতে হয়।\\\hline
বিক্রির CSV বা রেকর্ড & হিসাব, চার্ট বা পূর্বাভাসের সূত্র, একক ও নমুনা হিসাব মিলিয়ে দেখতে হয়।\\\hline
কোড ও ত্রুটির বার্তা & সংশোধন ও পরীক্ষার প্রস্তাব কোড চালিয়ে যাচাই করতে হয়।\\\hline
লক্ষ্য ও টুলের ফল & তৈরি ফাইল, প্রতিবেদন বা অনুমোদিত কাজের লগ ও ফল পরীক্ষা করতে হয়।\\\hline
\end{longtable}
\end{compacttable}
স্পষ্ট প্রম্পটে লক্ষ্য, প্রয়োজনীয় তথ্য, পাঠকের স্তর ও ফলের গঠন থাকে। যেমন, দেওয়া HTML-এ \texttt{src} ও \texttt{alt}-এর কাজ ব্যাখ্যা, ভুল পথের উদাহরণ এবং ফাইল না দেখে নিশ্চিত দাবি না করার নির্দেশ দেওয়া যায়। শুধু ভদ্রতা বা দীর্ঘ প্রম্পট নির্ভুলতা নিশ্চিত করে না।
